\section{Saturated decompositions and further consequences}\label{sec:consequences}

For a length-$r$ decomposition define the $A$-slot excess by
\[
 \Delta_A=\sum_{t=1}^r\rank A_t-27.
\]
The split bound shows that $\Delta_A\geq0$. A decomposition is
\emph{saturated in the $A$-slot} if $\Delta_A=0$.

\begin{corollary}[Zero terms and saturated profiles]\label{cor:profiles}
In an arbitrary length decomposition saturated in the $A$-slot,
including one with zero summands, at most one $A_t$ is invertible.
If its length is $22$ and all $A_t$ are nonzero, its $A$-slot
matrix-rank profile is either
\begin{equation}\label{eq:22profiles}
 (17,5,0)\quad\text{or}\quad(18,3,1).
\end{equation}
Without the nonzero-$A_t$ hypothesis, either these profiles hold or
deleting a zero first factor gives a saturated length-$21$
decomposition.
\end{corollary}
\begin{proof}
For arbitrary factors,
\[
 27\leq\sum_t\rank M(A_t,B_t,C_t)
       \leq\sum_t\rank A_t=27.
\]
Hence each individual rank deficit is zero. In particular if $A_t\ne0$,
neither $B_t$ nor $C_t$ can be zero. Removing all terms with $A_t=0$
leaves nonzero factors and preserves the rank sum, so
Proposition~\ref{prop:invertibility} applies.

For $22$ nonzero first factors, write $n_j$ for the count of rank $j$.
Then $n_1+n_2+n_3=22$ and $n_1+2n_2+3n_3=27$, whence
$n_2+2n_3=5$. Since $n_3\leq1$, the only possibilities
are~\eqref{eq:22profiles}. If a first factor is zero, deletion leaves
length $21$ and the same rank sum. Two such zero factors would produce
length $20$, contrary to Theorem~\ref{thm:main}.
\end{proof}

The cyclic symmetry
\[
 (A,B,C)\longmapsto(B,C^T,A^T)
\]
follows by relabelling $i,j,k$ in~\eqref{eq:tensor}, and its third
iterate is the identity. Transpose preserves matrix rank. Thus the
split bound, the invertibility obstruction, and
Corollary~\ref{cor:profiles} hold separately in the $B$- and $C$-slots.
Saturation in one slot does not imply saturation in either other slot.

\subsection{Positive rank excess}
At length $21$, the three occupation capacities driving the
length-$20$ argument change as follows:
\begin{center}
\begin{tabular}{lcc}
\toprule
Subspace family & Length $20$ & Length $21$\\
\midrule
Lines with bound $19$ & $1$ & $2$\\
All-high planes with bound $19$ & $1$ & $2$\\
Affine-plane spans with bound $17$ & $3$ & $4$\\
\bottomrule
\end{tabular}
\end{center}
Repeated first factors and two high-rank factors in an all-high plane
are no longer immediately excluded. Four coset labels forming an affine
plane are also allowed. The argument therefore no longer forces the
profile or rank equality. At positive excess the hypotheses of
Lemma~\ref{lem:saturation} are unavailable.

The excess has a precise linear-algebraic meaning. For nonzero
summands set $U_t=\im M_t$. The addition map
\[
 \Sigma:\bigoplus_t U_t\longrightarrow\F^{27},\qquad
 (u_t)_t\longmapsto\sum_t u_t
\]
is surjective because $\sum_tM_t=P$ is invertible. Rank--nullity gives
\begin{equation}\label{eq:excess-kernel}
 \dim\ker\Sigma=\sum_t\dim U_t-27=\Delta_A.
\end{equation}
Zero excess means that components are unique, which is exactly what
produces the identities in Lemma~\ref{lem:saturation}. Positive excess
measures the remaining linear relations between these components.
Extending the method to length $21$ therefore requires stronger
occupation bounds or cross-factor relations that remain effective
when this kernel has small positive dimension.
The interval $21$--$23$ leaves both improved constructions and stronger
lower bounds to be determined.

\subsection{Global frozen-table soundness}
The full catalogue extends the selected quotient bounds to a single
lower-bound function on all subspaces of $V$.

\begin{proposition}\label{prop:global}
For every frozen representative and every subspace $W\leq V$,
\[
 R_{\F}(T_{W_i})\geq\ell_i\quad(0\leq i<496),
 \qquad R_{\F}(T_W)\geq L_0(W).
\]
\end{proposition}
\begin{proof}[Computer-assisted proof]
The completed registry supplies a closed proof at each of the $496$
listed bases and labels. Its bindings check equal spans, containment
with quotient monotonicity, or invertible tensor transport to the
individual source proofs. The entry $W_{495}=0$, labelled $20$, may
now use Theorem~\ref{thm:main}, since $T_0=\T$. For each $G$-image,
tensor symmetry preserves the quotient rank. Taking the maximum in
\eqref{eq:L0} gives the second inequality.
\end{proof}

With universal soundness available, there is also a direct link from
any short quotient decomposition to the full occupation system. At
length $r\leq18$, its nonzero first-factor multiplicities satisfy the
caps $r-L_0(U)$. Increasing any one multiplicity by $18-r$ preserves
every enlarged cap $18-L_0(U)$ and makes the total $18$. This explains
the padding argument for the full systems used in the research report.
The proof above instead extracts the retained system from separately
established source bounds, as in the Lean main-theorem dependency graph.
